\section{\texorpdfstring{射影 $\mathcal{V}$ 算子}{射影 V 算子}}\label{sec:projective-vanishing-operator}

与仿射簇一样, 射影簇也定义为一族多项式的共同零点集. 按 Definition~\ref{def:projective-space-homogeneous-coordinates} 对 $\mathbb{P}^{n}$ 的描述, 这些多项式的输入是元组 $(a_{0},\ldots,a_{n})$. 然而, Definition~\ref{def:projective-space-homogeneous-coordinates} 涉及等价关系, 所以一个多项式可能在一个代表元上取值为零, 而在另一个代表元上不为零. 因此, 多项式在射影空间的一点处\textquotedbl{}取值为零\textquotedbl{}的含义并不显然.

例如, 将 $\mathbb{P}^{1}$ 中形如 $[a_{0}:a_{1}]$ 的点作为二元多项式 $f=x_{0}^{2}-x_{1}$ 的输入. 点 $[2:4]$ 看起来属于 $f$ 的零点集, 因为
\[
  f(2,4)=2^{2}-4=0.
\]
然而, 另一方面, $[2:4]=[4:8]$, 且
\[
  f(4,8)=4^{2}-8=8\neq0.
\]
因此, $f$ 是否在点 $[2:4]=[4:8]$ 处取值为零, 似乎没有一个良定义的答案. 解决这一不一致的方法是规定: 只有当一个多项式在射影空间某点的每个代表元上取值都为零时, 才说它在该点处\textquotedbl{}取值为零\textquotedbl{}.

\begin{Definition}{射影消没}\leavevmode\label{def:projective-vanishing}
设 $f\in K[x_{0},\ldots,x_{n}]$ 为多项式, $a\in\mathbb{P}^{n}$ 为一点. 若对每个代表元 $a=[a_{0}:\cdots:a_{n}]$ 都有
\[
  f(a_{0},\ldots,a_{n})=0,
\]
则称 $f$ \emph{在 $a$ 处取值为零}, 并记为 $f(a)=0$.
\end{Definition}

例如, 多项式 $f(x_{0},x_{1})=x_{0}^{2}-x_{1}$ 在点 $a=[2:4]\in\mathbb{P}^{1}$ 处不为零, 因为它在等价代表元 $a=[4:8]$ 上取值不为零. 乍看之下, 检查多项式是否在一个点的每个代表元上取值都为零, 似乎很费力. 毕竟, 每个点都有无穷多个代表元. 不过, 引入齐次多项式后, 这项工作可以简化.

\begin{Definition}{齐次多项式}\leavevmode\label{def:homogeneous-polynomial}
若多项式 $f\in K[x_{0},\ldots,x_{n}]$ 的每一非零项的次数都是 $d$, 则称 $f$ 是 \emph{$d$ 次齐次多项式}.
\end{Definition}

\begin{BookFlowNote}[l]{%
例如, 二元多项式 $f=x_{0}^{2}-x_{1}$ 是非齐次的, 因为它有一项是非零的二次项, 另一项是非零的一次项; 而 $g=x_{0}^{2}-2x_{0}x_{1}$ 是二次齐次多项式.
}
\noindent $f$ 的一\textquotedbl{}项\textquotedbl{}是表达式 $f=\sum_{\alpha}a_{\alpha}x^{\alpha}$ 中的任何一个加项. 当 $a_{\alpha}\neq0$ 时, 项 $a_{\alpha}x^{\alpha}$ 是\textquotedbl{}非零的\textquotedbl{}.

零多项式没有非零项, 所以它对任何次数都满足齐次性.
\end{BookFlowNote}

在研究射影空间中的消没时, 使用齐次多项式的重要性体现在下面的结果中.

\begin{Lemma}{齐次多项式的射影消没}\leavevmode\label{lem:homogeneous-polynomial-vanishing-on-representatives}
设 $f\in K[x_{0},\ldots,x_{n}]$ 为齐次多项式, 且
\[
  [a_{0}:\cdots:a_{n}]=[b_{0}:\cdots:b_{n}]\in\mathbb{P}^{n}.
\]
则
\[
  f(a_{0},\ldots,a_{n})=0\quad\Longleftrightarrow\quad f(b_{0},\ldots,b_{n})=0.
\]
\end{Lemma}

换句话说, 对齐次多项式, 只须检查它在一个代表元上的取值是否为零. 例如, 考虑齐次多项式 $g=x_{0}^{2}-2x_{0}x_{1}$. 要验证 $g$ 在等价类 $[2:1]\in\mathbb{P}^{1}$ 处取值为零, 只须验证一个代表元:
\[
  g(2,1)=2^{2}-2\cdot2\cdot1=0.
\]
若把 $[2:1]$ 换成另一个代表元 $[4:2]$ (或 $\mathbb{P}^{1}$ 中这个点的任一其他代表元), 取值仍为零:
\[
  g(4,2)=4^{2}-2\cdot4\cdot2=0.
\]

\begin{Proof}\textbf{of Lemma~\ref{lem:homogeneous-polynomial-vanishing-on-representatives}.}\quad
所需的关键观察是, $f$ 为 $d$ 次齐次多项式当且仅当对所有 $\lambda,a_{0},\ldots,a_{n}\in K$, 有
\begin{equation}\label{eq:homogeneous-polynomial-scaling}
  f(\lambda a_{0},\ldots,\lambda a_{n})=\lambda^{d}f(a_{0},\ldots,a_{n})
\end{equation}
(Exercise~\ref{ex:characterize-homogeneous-polynomials-by-scaling}). 设 $f$ 为 $d$ 次齐次多项式, 且
\[
  [a_{0}:\cdots:a_{n}]=[b_{0}:\cdots:b_{n}]\in\mathbb{P}^{n}.
\]
按 $\mathbb{P}^{n}$ 上等价关系的定义, 存在非零的 $\lambda\in K$, 使得
\[
  (\lambda a_{0},\ldots,\lambda a_{n})=(b_{0},\ldots,b_{n}).
\]
由 \eqref{eq:homogeneous-polynomial-scaling} 可见, $f(b_{0},\ldots,b_{n})$ 与 $f(a_{0},\ldots,a_{n})$ 相差非零的标量倍数 $\lambda^{d}$, 所以一个为零当且仅当另一个为零.
\end{Proof}

Lemma~\ref{lem:homogeneous-polynomial-vanishing-on-representatives} 表明, 只须检查一个代表元, 便可判定齐次多项式是否在射影空间的一点处取值为零. 但非齐次多项式又如何? 为回答这个问题, 我们引入多项式的齐次分量.

\begin{Definition}{齐次分量}\leavevmode\label{def:homogeneous-components}
给定多项式 $f\in K[x_{0},\ldots,x_{n}]$, $f$ 的 \emph{$d$ 次齐次分量}记为 $f_{d}$, 是 $f$ 的所有 $d$ 次非零项之和. 若 $f$ 没有 $d$ 次非零项, 则 $f_{d}=0$.
\end{Definition}

例如, 多项式
\[
  f=x^{4}z+x^{2}y+xyz+x+y+5
\]
的非零齐次分量是
\[
  f_{5}=x^{4}z,\quad f_{3}=x^{2}y+xyz,\quad f_{1}=x+y,\quad f_{0}=5.
\]
下一个结果用齐次分量的消没来刻画多项式在射影空间一点处的消没.

\begin{Lemma}{射影消没与齐次分量}\leavevmode\label{lem:projective-vanishing-of-components}
设 $f\in K[x_{0},\ldots,x_{n}]$ 是多项式, $a\in\mathbb{P}^{n}$ 为一点. 则 $f$ 在 $a$ 处取值为零当且仅当 $f$ 的每个齐次分量都在 $a$ 处取值为零.
\end{Lemma}

换句话说, 为判定一般的多项式是否在 $\mathbb{P}^{n}$ 中一点 (的每个代表元) 处取值为零, 由 Lemmas~\ref{lem:homogeneous-polynomial-vanishing-on-representatives} 和~\ref{lem:projective-vanishing-of-components}, 只须检查每个齐次分量在该点的一个代表元上是否都取值为零.

\begin{Proof}\textbf{of Lemma~\ref{lem:projective-vanishing-of-components}.}\quad
设 $f\in K[x_{0},\ldots,x_{n}]$ 是 $d$ 次多项式. 它可写成次数 $\leq d$ 的各齐次分量之和:
\begin{equation}\label{eq:polynomial-sum-of-homogeneous-components}
  f=\sum_{k=0}^{d}f_{k}.
\end{equation}
若每个 $f_{k}$ 在点 $a$ 的每个代表元上取值都为零, 则由 \eqref{eq:polynomial-sum-of-homogeneous-components}, $f$ 在 $a$ 的每个代表元上取值都为零.

反过来, 假设 $f$ 在 $a$ 的每个代表元上取值都为零. 选定一个代表元 $a=[a_{0}:\cdots:a_{n}]$. 那么对任意 $\lambda\in K\setminus\{0\}$, 有
\[
  \begin{aligned}
    0=f(\lambda a_{0},\ldots,\lambda a_{n})
      &=\sum_{k=0}^{d}f_{k}(\lambda a_{0},\ldots,\lambda a_{n})\\
      &=\sum_{k=0}^{d}\lambda^{k}f_{k}(a_{0},\ldots,a_{n}),
  \end{aligned}
\]
其中第二个等号用了 \eqref{eq:homogeneous-polynomial-scaling}. 换句话说, 一元多项式
\[
  \sum_{k=0}^{d}x^{k}f_{k}(a_{0},\ldots,a_{n})\in K[x]
\]
在 $K$ 中无穷多个值处取值为零, 故必为零多项式. 于是对所有 $k$, 都有 $f_{k}(a_{0},\ldots,a_{n})=0$. 由于 $f_{k}$ 齐次, Lemma~\ref{lem:homogeneous-polynomial-vanishing-on-representatives} 蕴含每个 $f_{k}$ 在 $a$ 的每个代表元上取值都为零.
\end{Proof}

更清楚地理解了多项式在射影空间的点处取值为零的含义之后, 现在给出射影簇的自然定义.

\begin{Definition}{射影 $\mathcal{V}$ 算子}\leavevmode\label{def:projective-vanishing-operator}
设 $S\subseteq K[x_{0},\ldots,x_{n}]$ 为多项式集合. $S$ 的\emph{射影零点集}是
\[
  \mathcal{V}_{\mathbb{P}}(S)=\{a\in\mathbb{P}^{n}\mid f(a)=0\text{ 对所有 }f\in S\}.
\]
若子集 $X\subseteq\mathbb{P}^{n}$ 对某个集合 $S\subseteq K[x_{0},\ldots,x_{n}]$ 满足 $X=\mathcal{V}_{\mathbb{P}}(S)$, 则称 $X$ 为\emph{射影簇}.
\end{Definition}

\begin{BookFlowNote}[l]{%
由 Lemma~\ref{lem:projective-vanishing-of-components}, 只须把 $S$ 中的非齐次多项式换成它们的齐次分量, 就能用一组齐次多项式描述每个射影簇.
}
\noindent 当由语境可清楚看出是在射影空间, 而非仿射空间中讨论时, 常写作 $\mathcal{V}(S)\subseteq\mathbb{P}^{n}$.

因此, 通常把射影簇描述为一族齐次多项式的零点集. 来看几个例子.
\end{BookFlowNote}

\begin{Example}{$\varnothing$ 和 $\mathbb{P}^{n}$ 是射影簇}\leavevmode\label{ex:empty-set-and-projective-space-are-projective-varieties}
与仿射情形一样, $\mathcal{V}(1)=\varnothing$ 且 $\mathcal{V}(0)=\mathbb{P}^{n}$, 所以 $\varnothing$ 和 $\mathbb{P}^{n}$ 都是射影簇.
\end{Example}

\begin{Example}{$\mathbb{P}^{1}$ 中的射影簇}\leavevmode\label{ex:projective-varieties-in-projective-line}
在 $\mathbb{P}^{1}$ 中, 考虑射影簇
\[
  X=\mathcal{V}(2x_{0}-x_{1})
   =\{[a_{0}:a_{1}]\in\mathbb{P}^{1}\mid 2a_{0}-a_{1}=0\}.
\]
$X$ 中的点 $[a_{0}:a_{1}]$ 不能满足 $a_{0}=0$, 否则方程 $2a_{0}-a_{1}=0$ 会迫使 $a_{1}$ 也为零. 因此
\[
  X=\{[a:2a]\in\mathbb{P}^{1}\mid a\in K\setminus\{0\}\}=\{[1:2]\},
\]
因为将两个坐标同时乘以 $a^{-1}$ 可知, 对任意 $a$, 有 $[a:2a]=[1:2]$. 更一般地, 与 $\mathbb{A}^{1}$ 的情形完全类似, $\mathbb{P}^{1}$ 中的任意射影簇要么是整个 $\mathbb{P}^{1}$, 要么是有限 (可能为空) 的点集 (Exercise~\ref{ex:classify-varieties-in-projective-line}).
\end{Example}

\begin{Example}{$\mathbb{P}^{2}$ 中的一条直线}\leavevmode\label{ex:projective-line-in-projective-plane}
在 $\mathbb{P}^{2}$ 中, 考虑射影簇
\[
  X=\mathcal{V}(x_{0}+x_{1}-x_{2})
   =\{[a_{0}:a_{1}:a_{2}]\in\mathbb{P}^{2}\mid a_{0}+a_{1}-a_{2}=0\}.
\]
$X$ 中的点 $[a_{0}:a_{1}:a_{2}]$ 不能满足 $a_{0}=a_{1}=0$, 否则定义方程会迫使 $a_{2}$ 也为零. 从而
\[
  X=\{[a_{0}:a_{1}:a_{0}+a_{1}]\in\mathbb{P}^{2}\mid a_{0},a_{1}\in K\text{ 不全为 }0\}.
\]
由此不难看出, $\mathcal{V}(x_{0}+x_{1}-x_{2})$ 的点与 $\mathbb{P}^{1}$ 有双射.

与仿射情形类似, 在 $\mathbb{A}^{2}$ 中线性多项式的零点集描述一条直线, 我们称 $X$ 为 $\mathbb{P}^{2}$ 中的一条直线. 这个直观的术语不宜完全照字面理解: 如上一节所见, 射影\textquotedbl{}直线\textquotedbl{}实际上不像平面几何中熟悉的直线. 即使在实数域上也是如此, 此时射影直线形成一个环.
\end{Example}

要利用多项式环的代数结构研究射影簇, 第一步是注意每个射影簇都可由一个理想定义. 这个结果与仿射情形中的 Proposition~\ref{prop:affine-varieties-defined-by-ideals} 对应.

\begin{Proposition}{射影簇可由理想定义}\leavevmode\label{prop:projective-varieties-defined-by-ideals}
若 $S\subseteq K[x_{0},\ldots,x_{n}]$ 为多项式集合, 则
\[
  \mathcal{V}_{\mathbb{P}}(S)=\mathcal{V}_{\mathbb{P}}(\langle S\rangle).
\]
\end{Proposition}
\begin{Proof}
见 Exercise~\ref{ex:projective-vanishing-of-generated-ideal}.
\end{Proof}

和仿射情形一样, 知道每个射影簇都可由一个理想定义后, 就能利用多项式环的代数结构推出, 每个射影簇都可由有限多个多项式定义. 而且, 在射影情形还能进一步利用 Lemma~\ref{lem:projective-vanishing-of-components} 推出, 每个射影簇都是一组有限个齐次多项式的零点集.

\begin{Corollary}{射影簇可由有限个多项式定义}\leavevmode\label{cor:projective-varieties-defined-by-finite-homogeneous-set}
任意射影簇 $X\subseteq\mathbb{P}^{n}$ 都具有形式
\[
  X=\mathcal{V}_{\mathbb{P}}(f_{1},\ldots,f_{k}),
\]
其中 $f_{1},\ldots,f_{k}\in K[x_{0},\ldots,x_{n}]$ 为齐次多项式.
\end{Corollary}
\begin{Proof}
见 Exercise~\ref{ex:finite-homogeneous-equations-for-projective-variety}.
\end{Proof}

到这里, 已可看出用射影簇的定义理想, 而非仅用多项式集合来作代数研究的好处. 特别地, 正是通过转到理想并应用多项式环的 Hilbert 基定理, 才能证明 Corollary~\ref{cor:projective-varieties-defined-by-finite-homogeneous-set}. 和仿射情形一样, 一个射影簇可以有许多定义理想, 但总有一个特别的定义理想, 即消没理想. 下一节转而讨论射影情形中的消没理想.

\BookExerciseHeading
\begin{BookExercises}
  \item\label{ex:characterize-homogeneous-polynomials-by-scaling} 设 $f\in K[x_{0},x_{1},\ldots,x_{n}]$ 为非零多项式. 证明 $f$ 为 $d$ 次齐次多项式当且仅当对任意 $\lambda,a_{0},a_{1},\ldots,a_{n}\in K$, 都有
  \[
    f(\lambda a_{0},\lambda a_{1},\ldots,\lambda a_{n})=\lambda^{d}f(a_{0},a_{1},\ldots,a_{n}).
  \]
  \item\label{ex:classify-varieties-in-projective-line} 证明 $\mathbb{P}^{1}$ 中的射影簇只有 $\mathbb{P}^{1}$ 本身和有限点集.
  \item\label{ex:projective-vanishing-of-generated-ideal} 证明 Proposition~\ref{prop:projective-varieties-defined-by-ideals}.
  \item\label{ex:finite-homogeneous-equations-for-projective-variety} 证明 Corollary~\ref{cor:projective-varieties-defined-by-finite-homogeneous-set}.
  \item\label{ex:projective-varieties-closed-under-unions-and-intersections} 证明射影簇的有限并和任意交都是射影簇.
  \item\label{ex:projective-conic-and-affine-chart} 令
  \[
    X=\mathcal{V}(x_{0}x_{1}-x_{2}^{2})\subseteq\mathbb{P}^{2}.
  \]
  \begin{enumerate}
    \item 令 $U=\{[a_{0}:a_{1}:a_{2}]\in\mathbb{P}^{2}\mid a_{0}\neq0\}$. 由上一节的结论, 它与 $\mathbb{A}^{2}$ 有自然双射. 证明这一双射将 $X\cap U$ 与 $\mathbb{A}^{2}$ 中的一个仿射簇认同. 这是哪个仿射簇? 画出 $X\cap U$ 在 $\mathbb{R}$ 上的图形.
    \item 求 $X\setminus U$ 的所有点, 并描述 (b) 中找到的仿射簇里的点如何趋近这些点.
    \item 描述 $X$ 与 $\mathbb{P}^{1}$ 之间的一个双射.
  \end{enumerate}
  \item\label{ex:projective-conic-with-other-affine-chart} 对 $X=\mathcal{V}(x_{0}^{2}-x_{1}x_{2})\subseteq\mathbb{P}^{2}$ 完成上一题.
\end{BookExercises}

\par\medskip
{\footnotesize
\noindent\textbf{开放获取.} 本章按照
\href{http://creativecommons.org/licenses/by-nc/4.0/}{Creative Commons Attribution-NonCommercial 4.0 International License}
的条款授权. 该许可允许以任何媒介或格式进行非商业性使用, 分享, 改编, 分发和复制, 但须对原作者和来源作适当署名, 提供 Creative Commons 许可的链接, 并说明是否作了修改.

\noindent
\begin{minipage}[t]{.17\linewidth}
  \vspace{0pt}\includegraphics[width=.9\linewidth]{figures/source-cc-by-nc.png}
\end{minipage}\hfill
\begin{minipage}[t]{.81\linewidth}
  \vspace{0pt}本章的图片或其他第三方材料均纳入本章的 Creative Commons 许可, 除非材料的署名行另有注明. 如果材料未纳入本章的 Creative Commons 许可, 且您拟进行的使用未获法律规定的许可或超出允许使用的范围, 则须直接向著作权持有人取得许可.
\end{minipage}\par}
